\documentclass[11pt,aspectratio=169]{beamer}
\usepackage[italian]{babel}
\usepackage[latin1]{inputenc}
\usepackage{graphicx}
\usepackage{listings}
\usepackage[export]{adjustbox}
% Custom bullets
\usepackage{pifont}
\usepackage[useregional]{datetime2}
\usepackage{tikz}
\usepackage[table]{colortbl}
\usepackage{venndiagram}
\usepackage{amsmath}

% custom legend
\newcommand{\cbox}[1]{\raisebox{\depth}{\fcolorbox{black}{#1}{\null}}}

\usetheme{Madrid}
\usecolortheme{spruce}

% Redefine labels
\deftranslation[to=Italian]{Section}{Sezione}
\deftranslation[to=Italian]{Subsection}{Sottosezione}

% Join commands
\def\ojoin{\setbox0=\hbox{$\bowtie$}%
  \rule[-.02ex]{.25em}{.4pt}\llap{\rule[\ht0]{.25em}{.4pt}}}
\def\leftouterjoin{\mathbin{\ojoin\mkern-5.8mu\bowtie}}
\def\rightouterjoin{\mathbin{\bowtie\mkern-5.8mu\ojoin}}
\def\fullouterjoin{\mathbin{\ojoin\mkern-5.8mu\bowtie\mkern-5.8mu\ojoin}}

\author[Piero Hierro \and Piermichele Rosati]{Piero Hierro \and Piermichele Rosati}

\institute[]{\large Universit\`a di Camerino\\ \footnotesize Tutorato - Basi di Dati}

\title[Esercizi - Algebra Relazionale e Fondamenti SQL]{3. Esercizi - Algebra Relazionale e Fondamenti SQL}
\setbeamertemplate{navigation symbols}{}
\setbeamertemplate{section in toc}[sections numbered]
\setbeamertemplate{subsection in toc}[subsections numbered]
%\titlegraphic{\includegraphics[width=6cm]{img/unicam-logo.jpg}}
\makeatletter
\setbeamertemplate{footline}{
    \leavevmode%
    \hbox{%
        \begin{beamercolorbox}[wd=.333333\paperwidth,ht=2.25ex,dp=1ex,center]{author in head/foot}%
            \usebeamerfont{author in head/foot}\insertshortauthor
        \end{beamercolorbox}%
        \begin{beamercolorbox}[wd=.333333\paperwidth,ht=2.25ex,dp=1ex,center]{title in head/foot}%
            \usebeamerfont{title in head/foot}\insertshorttitle
        \end{beamercolorbox}%
        \begin{beamercolorbox}[wd=.333333\paperwidth,ht=2.25ex,dp=1ex,right]{date in head/foot}%
            \usebeamerfont{date in head/foot}\insertshortdate\hspace*{5 em}
            \insertframenumber{} / \inserttotalframenumber\hspace*{2ex} 
        \end{beamercolorbox}%
    }%
    \vskip0pt%
}
\makeatother


\AtBeginSection[]
{
\begin{frame}<beamer>
\frametitle{Indice}
\tableofcontents[currentsection]
\end{frame}
}

\begin{document}
\begin{frame}
\centering
\includegraphics[width=5.5cm]{../img/unicam-logo.jpg}
%\date{17 Aprile 2024}
\titlepage
\end{frame}
\addtobeamertemplate{frametitle}{}{%
\begin{tikzpicture}[remember picture,overlay]
\node[anchor=north east,yshift=2pt] at (current page.north east) {\includegraphics[height=2cm]{../img/unicam-logo-notext.png}};
\end{tikzpicture}}

%\begin{frame}
%\titlepage
%\end{frame}

\section{Algebra Relazionale}
\def\schemaEx3.8{\small \begin{align*}
    DEPUTATI &= (\underline{Codice}, Cognome, Nome, Commissione, Provincia, Collegio)\\
    COLLEGI &= (\underline{Provincia}, Numero,Nome)\\
    PROVINCE &= (\underline{Sigla}, Nome, Regione)\\
    REGIONI &= (\underline{Codice}, Nome)\\
    COMMISSIONI &= (\underline{Numero}, Nome, Presidente)\\
    \end{align*}}
\begin{frame}{Esercizio 3.8}
    \framesubtitle{3.8.1}
    \schemaEx3.8
    Trovare nome e cognome dei presidenti di commissioni cui partecipa almeno un deputato eletto in una provincia della Sicilia.
\end{frame}
%
\begin{frame}{Esercizio 3.8}
    \framesubtitle{Soluzione 3.8.1}
    \vspace*{-1.2cm}
    \schemaEx3.8
    {\small Nome e cognome dei presidenti di commissioni cui partecipa almeno un deputato eletto in una provincia della Sicilia:}
    \small
    \begin{gather*}
        \onslide<8->{\pi_{Nom,Cogn}}\\
        \onslide<7->{((\rho_{Nom,Cogn~\leftarrow~Nome,Cognome}}
        \onslide<6->{(DEPUTATI))~\bowtie_{Presidente~=~Codice}}\\
        \onslide<5->{(COMMISSIONI~\bowtie_{Numero~=~Comm}}
        \onslide<4->{(\rho_{Comm~\leftarrow~{Commissione}}}
        \onslide<3->{(DEPUTATI~\bowtie_{Provincia~=~Sigla}}\\
        \onslide<2->{(PROVINCE~\bowtie_{Regione~=~Codice}}\\
        \onslide<1->{(\sigma_{Nome~=~'Sicilia'}(REGIONI))}
        \onslide<2->{)}
        \onslide<3->{)}
        \onslide<4->{)}
        \onslide<5->{)}
        \onslide<6->{)}
    \end{gather*}
\end{frame}
%
\begin{frame}{Esercizio 3.8}
    \framesubtitle{3.8.2}
    \schemaEx3.8
    Trovare nome e cognome dei deputati della commissione Bilancio.
\end{frame}
%
\begin{frame}{Esercizio 3.8}
    \framesubtitle{Soluzione 3.8.2}
    \vspace*{-1.2cm}
    \schemaEx3.8
    {\small Nome e cognome dei deputati della commissione Bilancio:}
    \small
    \begin{gather*}
        \onslide<4->{\pi_{NomeC,Cognome}(}
        \onslide<3->{(\rho_{NomeC~\leftarrow~Nome}}
        \onslide<2->{(DEPUTATI))~\bowtie_{Commissione~=~Numero}}\\
        \onslide<1->{(\sigma_{Nome~=~'Bilancio'}(COMMISSIONI))}
    \end{gather*}
\end{frame}
%
\begin{frame}{Esercizio 3.8}
    \framesubtitle{3.8.3}
    \schemaEx3.8
    Trovare nome, cognome e provincia di elezione dei deputati della commissione Bilancio.
\end{frame}
%
\begin{frame}{Esercizio 3.8}
    \framesubtitle{Soluzione 3.8.3}
    \vspace*{-1.2cm}
    \schemaEx3.8
    {\small Nome, cognome e provincia di elezione dei deputati della commissione Bilancio:}
    \small
    \begin{gather*}
        \onslide<6->{\pi_{NomeC,Cognome,NomeP}(}\\
        \onslide<5->{(\rho_{NomeP~\leftarrow~Nome}}
        \onslide<4->{(PROVINCE))~\bowtie_{Sigla~=~Provincia}}\\
        \onslide<3->{((\rho_{NomeC~\leftarrow~Nome}}
        \onslide<2->{(DEPUTATI))~\bowtie_{Commissione~=~Numero}}\\
        \onslide<1->{(\sigma_{Nome~=~'Bilancio'}(COMMISSIONI))}
        \onslide<3->{)}
    \end{gather*}
\end{frame}
%
\begin{frame}{Esercizio 3.8}
    \framesubtitle{3.8.4}
    \schemaEx3.8
    Trovare nome, cognome, provincia e regione di elezione dei deputati della commissione Bilancio.
\end{frame}
%
\begin{frame}{Esercizio 3.8}
    \framesubtitle{Soluzione 3.8.4}
    \vspace*{-1.2cm}
    \schemaEx3.8
    {\small Nome, cognome, provincia e regione di elezione dei deputati della commissione Bilancio:}
    \small
    \begin{gather*}
        \onslide<8->{\pi_{NomeC,Cognome,NomeP,NomeR}(}\\
        \onslide<7->{(\rho_{NomeR~\leftarrow~Nome}}
        \onslide<6->{(REGIONI))~\bowtie_{Codice~=~Regione}}\\
        \onslide<5->{(\rho_{NomeP~\leftarrow~Nome}}
        \onslide<4->{(PROVINCE))~\bowtie_{Sigla~=~Provincia}}\\
        \onslide<3->{((\rho_{NomeC~\leftarrow~Nome}}
        \onslide<2->{(DEPUTATI))~\bowtie_{Commissione~=~Numero}}\\
        \onslide<1->{(\sigma_{Nome~=~'Bilancio'}(COMMISSIONI))}
        \onslide<3->{)}
    \end{gather*}
\end{frame}
%
\section{Fondamenti SQL}
%\begin{frame}{QUERY}
    \begin{table}[h]
    \centering
    \begin{tabular}{|c|c|c|c|}
    \hline
    \rowcolor{cyan!30}\multicolumn{4}{|c|}{Table} \\
    \hline
    \rowcolor{cyan!30}Column1 & Column2 & Column3 & Column4 \\
    \hline
    Data1 & Data2 & Data3 & Data4 \\
    Data5 & Data6 & Data7 & Data8 \\
    Data9 & Data10 & Data11 & Data12 \\
    \hline
    \end{tabular}
    \begin{tikzpicture}[overlay,remember picture]
      \draw[->,thick,line width=2pt] (1.5,-0.5) -- (1.5,-1.5) node[midway,right] {Query};
    \end{tikzpicture}
    \end{table}
    
    \begin{table}[h]
    \centering
    \begin{tabular}{|c|c|c|}
    \hline
    \rowcolor{cyan!30}\multicolumn{3}{|c|}{Temp Table} \\
    \hline
    \rowcolor{cyan!30}TempColumn1 & TempColumn2 & TempColumn3 \\
    \hline
    TempData1 & TempData2 & TempData3 \\
    TempData4 & TempData5 & TempData6 \\
    TempData7 & TempData8 & TempData9 \\
    \hline
    \end{tabular}
    \end{table}
    \end{frame}
    
    
    %%%%%%%%%%%%%%%%%%%%%%%%%%%
    \begin{frame}{SELECT}
    \begin{table}[h]
    \centering
    \begin{tabular}{|c|c|c|c|}
    \hline
    \rowcolor{cyan!30}\multicolumn{4}{|c|}{Table} \\
    \hline
    \rowcolor{cyan!30}column1 & column2 & column3 & column4 \\
    \hline
    data1 & data2 & data3 & data4 \\
    data5 & data6 & data7 & data8 \\
    data9 & data10 & data11 & data12 \\
    \hline
    \end{tabular}
    \end{table}
    
    \vspace{2em} 
    
    \texttt{SELECT column1, column2 \\FROM TABLE}
    \end{frame}
    %%%%%%%%%%%%%%%%%%%%%
    \begin{frame}{SELECT}
    \begin{table}[h]
    \centering
    \begin{tabular}{|c|c|c|c|}
    \hline
    \rowcolor{cyan!30}\multicolumn{4}{|c|}{Table} \\
    \hline
    \rowcolor{cyan!30}column1 & column2 & column3 & column4 \\
    \hline
    \cellcolor{red!20} data1 & \cellcolor{red!20} data2 & data3 & data4 \\
    \cellcolor{red!20} data5 & \cellcolor{red!20} data6 & data7 & data8 \\
    \cellcolor{red!20} data9 & \cellcolor{red!20} data10 & data11 & data12 \\
    \hline
    \end{tabular}
    \end{table}
    \vspace{2em} 
    
    \texttt{SELECT column1, column2 \\FROM TABLE}
    \end{frame}
    
    %%%%%%%%
    
    \begin{frame}{SELECT}
    \begin{table}[h]
    \centering
    \begin{tabular}{|c|c|c|c|}
    \hline
    \rowcolor{cyan!30}\multicolumn{4}{|c|}{Table} \\
    \hline
    \rowcolor{cyan!30}column1 &  column2 & column3 & column4 \\
    \hline
     data1 &  data2 & data3 & data4 \\
     data5 &  data6 & data7 & data8 \\
     data9 &  data10 & data11 & data12 \\
    \hline
    \end{tabular}
    \end{table}
    \vspace{2em} 
    
    \texttt{SELECT * \\FROM TABLE}
    \end{frame}
    
    %%%%%%%%%%%%%%
    \begin{frame}{SELECT}
    \begin{table}[h]
    \centering
    \begin{tabular}{|c|c|c|c|}
    \hline
    \rowcolor{cyan!30}\multicolumn{4}{|c|}{Table} \\
    \hline
    \rowcolor{cyan!30}column1 & column2 & column3 & column4 \\
    \hline
    \cellcolor{red!20}data1 & \cellcolor{red!20}data2 & \cellcolor{red!20}data3 & \cellcolor{red!20}data4 \\
    \cellcolor{red!20}data5 & \cellcolor{red!20}data6 & \cellcolor{red!20}data7 & \cellcolor{red!20}data8 \\
    \cellcolor{red!20}data9 & \cellcolor{red!20}data10 & \cellcolor{red!20}data11 & \cellcolor{red!20}data12 \\
    \hline
    \end{tabular}
    \end{table}
    
    \vspace{2em} 
    
    \texttt{SELECT * \\FROM TABLE}
    \end{frame}
    %%%%%%%%%%%%%%%
    \begin{frame}{SELECT}
    \begin{table}[h]
    \centering
    \begin{tabular}{|c|c|c|c|}
    \hline
    \rowcolor{cyan!30}\multicolumn{4}{|c|}{Table} \\
    \hline
    \rowcolor{cyan!30}nome &  column2 & column3 & column4 \\
    \hline
     Marco &  data2 & data3 & data4 \\
     Lucia &  data6 & data7 & data8 \\
     Marco &  data10 & data11 & data12 \\
    \hline
    \end{tabular}
    \end{table}
    \vspace{2em} 
    
    \texttt{SELECT nome \\FROM Table}
    \end{frame}
    %%%%%%%%%%%%%
    \begin{frame}{SELECT}
    \begin{table}[h]
    \centering
    \begin{tabular}{|c|c|c|c|}
    \hline
    \rowcolor{cyan!30}\multicolumn{4}{|c|}{Table} \\
    \hline
    \rowcolor{cyan!30}nome &  column2 & column3 & column4 \\
    \hline
    \cellcolor{red!20}Marco &  data2 & data3 & data4 \\
    \cellcolor{red!20}Lucia &  data6 & data7 & data8 \\
    \cellcolor{red!20}Marco &  data10 & data11 & data12 \\
    \hline
    \end{tabular}
    \begin{tikzpicture}[overlay,remember picture]
      \draw[->,thick,line width=2pt] (1.5,-0.5) -- (1.5,-1.5) node[midway,right] {Query};
    \end{tikzpicture}
    \end{table}
    
    \vspace{2em}
    
    \begin{table}[h]
    \centering
    \begin{tabular}{|c|}
    \hline
    \rowcolor{cyan!30}Result \\
    \hline
    Marco \\
    Lucia \\
    Marco \\
    \hline
    \end{tabular}
    \end{table}
    \end{frame}
    %%%%%%%%%
    \begin{frame}{SELECT}
    \begin{table}[h]
    \centering
    \begin{tabular}{|c|}
    \hline
    \rowcolor{cyan!30}Result \\
    \hline
    Marco \\
    Lucia \\
    \hline
    \end{tabular}
    \end{table}
    \vspace{2em}
    \texttt{SELECT DISTINCT nome \\FROM Table}
    \end{frame}
    %%%%%%%%
    
    \begin{frame}{Funzioni di aggregazione}
    \begin{table}[h]
    \centering
    \begin{tabular}{|c|c|}
    \hline
    \rowcolor{cyan!30}\multicolumn{2}{|c|}{Esami} \\
    \hline
    \rowcolor{cyan!30}voti &  ...  \\
    \hline
    18 &  ... \\
    30 &  ...  \\
    25 &  ...  \\
    27 &  ...  \\
    \hline
    \end{tabular}
    \end{table}
    
    \texttt{SELECT MIN(voti)}  \\
    \texttt{SELECT MAX(voti)}  \\
    \texttt{SELECT SUM(voti)}  \\
    \texttt{SELECT AVG(voti)}  \\
    \texttt{SELECT COUNT(voti)} 
    \end{frame}
    %%%%%%%%
    \begin{frame}{Funzioni di aggregazione}
    \begin{table}[h]
    \centering
    \begin{tabular}{|c|c|}
    \hline
    \rowcolor{cyan!30}\multicolumn{2}{|c|}{Esami} \\
    \rowcolor{cyan!30}\hline
    voti &  ...  \\
    \hline
    18 &  ... \\
    30 &  ...  \\
    25 &  ...  \\
    27 &  ...  \\
    \hline
    \end{tabular}
    \end{table}
    
    \texttt{SELECT MIN(voti)} \textcolor{red}{18} \\
    \texttt{SELECT MAX(voti)} \textcolor{red}{30} \\
    \texttt{SELECT SUM(voti)} \textcolor{red}{100} \\
    \texttt{SELECT AVG(voti)} \textcolor{red}{25} \\
    \texttt{SELECT COUNT(voti)} \textcolor{red}{4}
    \end{frame}
    %%%%%%%%
    \begin{frame}{FROM}
    L'istruzione SQL \texttt{FROM} viene utilizzata per specificare la tabella (o tabelle) da cui si desidera estrarre i dati.
    \end{frame}
    %%%%%
    \begin{frame}{WHERE}
    \begin{table}[h]
    \centering
    \begin{tabular}{|c|c|c|c|}
    \hline
    \rowcolor{cyan!30}\multicolumn{4}{|c|}{Table} \\
    \hline
    \rowcolor{cyan!30}column1 &  column2 & column3 & column4 \\
    \hline
     data1 &  data2 & data3 & data4 \\
     data5 &  data6 & data7 & data8 \\
     data9 &  data10 & data11 & data12 \\
    \hline
    \end{tabular}
    \end{table}
    \vspace{2em} 
    
    \texttt{SELECT * \\FROM TABLE\\ WHERE (condizione)}
    \end{frame}
    %%%%%
    \begin{frame}{WHERE}
    \begin{table}[h]
    \centering
    \begin{tabular}{|c|c|c|c|}
    \hline
    \rowcolor{cyan!30}\multicolumn{4}{|c|}{Table} \\
    \hline
    \rowcolor{cyan!30}column1 &  column2 & column3 & column4 \\
    \hline
     \cellcolor{red!20}data1 &  \cellcolor{red!20}data2 & \cellcolor{red!20}data3 & \cellcolor{red!20}data4 \\
     data5 &  data6 & data7 & data8 \\
     \cellcolor{red!20}data9 &  \cellcolor{red!20}data10 & \cellcolor{red!20}data11 & \cellcolor{red!20}data12 \\
    \hline
    \end{tabular}
    \end{table}
    \vspace{2em} 
    
    \texttt{SELECT * \\FROM TABLE\\ WHERE (condizione)}
    \end{frame}
    
    %%%%%%
    \begin{frame}{WHERE}
    \begin{table}[h]
    \centering
    \begin{tabular}{|c|c|}
    \hline
    \rowcolor{cyan!30}\multicolumn{2}{|c|}{Esami} \\
    \hline
    \rowcolor{cyan!30}voti &  ...  \\
    \hline
    18 &  ... \\
    30 &  ...  \\
    25 &  ...  \\
    27 &  ...  \\
    \hline
    \end{tabular}
    \end{table}
    \texttt{SELECT voti \\FROM Esami\\ WHERE voti>25;}
    \end{frame}
    %%%%%%%%%%%%%%
    \begin{frame}{WHERE}
    \begin{table}[h]
    \centering
    \begin{tabular}{|c|c|}
    \hline
    \rowcolor{cyan!30}\multicolumn{2}{|c|}{Esami} \\
    \hline
    \rowcolor{cyan!30}voti &  ...  \\
    \hline
    18 &  ... \\
    \cellcolor{red!20}{30} &  ...  \\
    25 &  ...  \\
    \cellcolor{red!20}{27} &  ...  \\
    \hline
    \end{tabular}
    \end{table}
    \texttt{SELECT voti \\FROM Esami\\ WHERE voti>25;}
    \end{frame}
    %%%%%%%%%%%
    \begin{frame}{WHERE}
    \begin{itemize}
        \item \textbf{Operatori di confronto:} $=, <>, >, <, >=, <=$
        \item \textbf{Operazioni booleane:} AND, OR, NOT
        \item \textbf{BETWEEN}
        \item \textbf{IN}
    \end{itemize}
    \end{frame}
    %%%%%%%%%%%
    \begin{frame}{WHERE}
    \begin{itemize}
        \item \textbf{Operatori di confronto:} $=, <>, >, <, >=, <=$
        \item \textbf{Operazioni booleane:} AND, OR, NOT
        \item \textbf{\textcolor{red}{BETWEEN}}
        \item \textbf{IN}
    \end{itemize}
    \vspace{2em}
    \texttt{WHERE Price BETWEEN 10 AND 20\\WHERE Price >= 10 AND Price <= 20}
    \end{frame}
    %%%%%%%%%%%%
    \begin{frame}{WHERE}
    \begin{itemize}
        \item \textbf{Operatori di confronto:} $=, <>, >, <, >=, <=$
        \item \textbf{Operazioni booleane:} AND, OR, NOT
        \item \textbf{BETWEEN}
        \item \textbf{IN}
    \end{itemize}
    \end{frame}
    %%%%%%%%%%%
    \begin{frame}{WHERE}
    \begin{itemize}
        \item \textbf{Operatori di confronto:} $=, <>, >, <, >=, <=$
        \item \textbf{Operazioni booleane:} AND, OR, NOT
        \item \textbf{BETWEEN}
        \item \textbf{\textcolor{red}{IN}}
    \end{itemize}
    \vspace{2em}
    \texttt{WHERE Country IN ('Germany', 'France', 'UK')\\WHERE Country = 'Germany' OR Country = 'France' OR Country = 'UK'}
    \end{frame}
    %%%%%%%%%%%%
    \begin{frame}{WHERE}
    \begin{itemize}
        \item \textbf{Operatori di confronto:} $=, <>, >, <, >=, <=$
        \item \textbf{Operazioni booleane:} AND, OR, NOT
        \item \textbf{BETWEEN}
        \item \textbf{IN}
        \item \textbf{\textcolor{red}{LIKE\{ \%, \_ \}}}
    \end{itemize}
    \vspace{2em}
    \textbf{\%} rappresenta zero, uno o pi\`u caratteri\\
    \textbf{\_} rappresenta un singolo carattere\\
    \end{frame}
    %%%%%%%%%%%%
    \begin{frame}{WHERE}
    \begin{itemize}
        \item \textbf{Operatori di confronto:} $=, <>, >, <, >=, <=$
        \item \textbf{Operazioni booleane:} AND, OR, NOT
        \item \textbf{BETWEEN}
        \item \textbf{IN}
        \item \textbf{\textcolor{red}{LIKE\{ \%, \_ \}}}
    \end{itemize}
    \vspace{2em}
    \texttt{SELECT * \\FROM Cliente\\WHERE nome LIKE 'A\%'}
    \end{frame}
    %%%%%%%%%%%%
    \begin{frame}{WHERE}
    \begin{columns}[T,onlytextwidth]
    \column{0.6\textwidth}
    \begin{itemize}
        \item \textbf{Operatori di confronto:} $=, <>, >, <, >=, <=$
        \item \textbf{Operazioni booleane:} AND, OR, NOT
        \item \textbf{BETWEEN}
        \item \textbf{IN}
        \item \textbf{\textcolor{red}{LIKE\{ \%, \_ \}}}
    \end{itemize}
    \vspace{2em}
    \texttt{SELECT * \\FROM Cliente\\WHERE nome LIKE 'A\%'}
    \column{0.5\textwidth}
    \centering
    \begin{table}[h]
    \centering
    \begin{tabular}{|c|c|}
    \hline
    \rowcolor{cyan!30}\multicolumn{2}{|c|}{Cliente} \\
    \hline
    \rowcolor{cyan!30}nome &  ...  \\
    \hline
    Anna &  ... \\
    Fabrizio &  ...  \\
    Asia &  ...  \\
    Laura &  ...  \\
    \hline
    \end{tabular}
    \end{table}
    \end{columns}
    \end{frame}
    %%%%%%%%%%%%
    \begin{frame}{WHERE}
    \begin{columns}[T,onlytextwidth]
    \column{0.6\textwidth}
    \begin{itemize}
        \item \textbf{Operatori di confronto:} $=, <>, >, <, >=, <=$
        \item \textbf{Operazioni booleane:} AND, OR, NOT
        \item \textbf{BETWEEN}
        \item \textbf{IN}
        \item \textbf{\textcolor{red}{LIKE\{ \%, \_ \}}}
    \end{itemize}
    \vspace{2em}
    \texttt{SELECT * \\FROM Cliente\\WHERE nome LIKE 'A\%'}
    \column{0.5\textwidth}
    \centering
    \begin{table}[h]
    \centering
    \begin{tabular}{|c|c|}
    \hline
    \rowcolor{cyan!30}\multicolumn{2}{|c|}{Cliente} \\
    \hline
    \rowcolor{cyan!30}nome &  ...  \\
    \hline
    \cellcolor{red!20}{Anna} &  ... \\
    Paolo &  ...  \\
    \cellcolor{red!20}{Asia} &  ...  \\
    Paola &  ...  \\
    \hline
    \end{tabular}
    \end{table}
    \end{columns}
    \end{frame}
    %%%%%%%%%%%%
    \begin{frame}{WHERE}
    \begin{columns}[T,onlytextwidth]
    \column{0.5\textwidth}
    \begin{itemize}
        \item \textbf{Operatori di confronto:} $=, <>, >, <, >=, <=$
        \item \textbf{Operazioni booleane:} AND, OR, NOT
        \item \textbf{BETWEEN}
        \item \textbf{IN}
        \item \textbf{\textcolor{red}{LIKE\{ \%, \_ \}}}
    \end{itemize}
    \vspace{2em}
    \texttt{SELECT * \\FROM Cliente\\WHERE nome LIKE 'Paol\_'}
    \column{0.6\textwidth}
    \centering
    \begin{table}[h]
    \centering
    \begin{tabular}{|c|c|}
    \hline
    \rowcolor{cyan!30}\multicolumn{2}{|c|}{Cliente} \\
    \hline
    \rowcolor{cyan!30}nome &  ...  \\
    \hline
    Anna &  ... \\
    Paolo &  ...  \\
    Asia &  ...  \\
    Paola &  ...  \\
    \hline
    \end{tabular}
    \end{table}
    \end{columns}
    \end{frame}
    %%%%%%%%%%%%
    \begin{frame}{WHERE}
    \begin{columns}[T,onlytextwidth]
    \column{0.6\textwidth}
    \begin{itemize}
        \item \textbf{Operatori di confronto:} $=, <>, >, <, >=, <=$
        \item \textbf{Operazioni booleane:} AND, OR, NOT
        \item \textbf{BETWEEN}
        \item \textbf{IN}
        \item \textbf{\textcolor{red}{LIKE\{ \%, \_ \}}}
    \end{itemize}
    \vspace{2em}
    \texttt{SELECT * \\FROM Cliente\\WHERE nome LIKE 'Paol\_'}
    \column{0.5\textwidth}
    \centering
    \begin{table}[h]
    \centering
    \begin{tabular}{|c|c|}
    \hline
    \rowcolor{cyan!30}\multicolumn{2}{|c|}{Cliente} \\
    \hline
    \rowcolor{cyan!30}nome &  ...  \\
    \hline
    Anna &  ... \\
    \cellcolor{red!20}{Paolo} &  ...  \\
    Asia &  ...  \\
    \cellcolor{red!20}{Paola} &  ...  \\
    \hline
    \end{tabular}
    \end{table}
    \end{columns}
    \end{frame}   
%\begin{frame}{INNER JOIN}
    \begin{center}
\begin{venndiagram2sets}[shade=red,tikzoptions={opacity=0.5},labelA={},labelB={}, showframe=false, ]
    \fillACapB
    \setpostvennhook
    {
        \draw (labelA) node[above=2pt] {A};
        \draw (labelB) node[above=2pt] {B};
    }
\end{venndiagram2sets}
\end{center}
\end{frame}
%%%%%%%%%%%%%%%%%
\begin{frame}{INNER JOIN}
    
\begin{table}[h]
\centering
\begin{minipage}{.45\textwidth}
\centering
\begin{tabular}{|c|c|c|}
\hline
\rowcolor{cyan!30}\multicolumn{3}{|c|}{Studente} \\
\hline
\rowcolor{cyan!30} matricola & nome & cognome \\
\hline
100 & Gianni & Morandi \\
101 & Jerry & Cal\`a \\
102 & Pippo & Franco \\
103 & Mara & Venier \\
\hline
\end{tabular}
\end{minipage}%
\begin{minipage}{.45\textwidth}
\centering
\begin{tabular}{|c|c|c|}
\hline
\rowcolor{cyan!30} \multicolumn{3}{|c|}{Esame} \\
\hline
\rowcolor{cyan!30}matricola & voto & corso \\
\hline
100 & 30 & 27035 \\
100 & 30 & 27038 \\
100 & 30 & 27010 \\
102 & 21 & 27038 \\
90 & 25 & 27045 \\
\hline
\end{tabular}
\end{minipage}
\end{table}
\vspace{2em}
\texttt{SELECT st.matricola, cognome, voto, corso\\FROM Studente AS st INNER JOIN Esame AS es ON st.matricola=es.matricola;}
\end{frame}
%%%%%%%%%%%%%%%%%
\begin{frame}{INNER JOIN}
\begin{table}[h]
\centering
\begin{tabular}{|c|c|c|c|}
\hline
\rowcolor{cyan!30} \multicolumn{4}{|c|}{Result} \\
\hline
\rowcolor{cyan!30} matricola & cognome & voto & corso \\
\hline
100 & Morandi & 30 & 27035 \\
100 & Morandi & 30 & 27038 \\
100 & Morandi & 30 & 27010 \\
102 & Franco & 21 & 27038 \\
\hline
\end{tabular}
\end{table}
\end{frame}
%%%%%%%%%%%%%%%%%
\begin{frame}{INNER JOIN}
    
\begin{table}[h]
\centering
\begin{minipage}{.45\textwidth}
\centering
\begin{tabular}{|c|c|c|}
\hline
\rowcolor{cyan!30}\multicolumn{3}{|c|}{Studente} \\
\hline
\rowcolor{cyan!30} matricola & nome & cognome \\
\hline
100 & Gianni & Morandi \\
101 & Jerry & Cal\`a \\
102 & Pippo & Franco \\
103 & Mara & Venier \\
\hline
\end{tabular}
\end{minipage}%
\begin{minipage}{.45\textwidth}
\centering
\begin{tabular}{|c|c|c|}
\hline
\rowcolor{cyan!30} \multicolumn{3}{|c|}{Esame} \\
\hline
\rowcolor{cyan!30}matricola & voto & corso \\
\hline
100 & 30 & 27035 \\
100 & 30 & 27038 \\
100 & 30 & 27010 \\
102 & 21 & 27038 \\
90 & 25 & 27045 \\
\hline
\end{tabular}
\end{minipage}
\end{table}
\vspace{2em}
\texttt{SELECT Studente.matricola, cognome, voto, corso\\ FROM Studente, Esame\\WHERE Studente.matricola=Esame.matricola;}
\end{frame}
%%%%%%%%%%%%%%%%%
\begin{frame}{LEFT JOIN}
\begin{center}
\begin{venndiagram2sets}[shade=red,tikzoptions={opacity=0.5},labelA={},labelB={}, showframe=false, ]
    \fillA
    \setpostvennhook
    {
        \draw (labelA) node[above=2pt] {A};
        \draw (labelB) node[above=2pt] {B};
    }
\end{venndiagram2sets}    
    \end{center}
\end{frame}
%%%%%%%%%%%%%%%%%
\begin{frame}{LEFT JOIN}
        
\begin{table}[h]
\centering
\begin{minipage}{.45\textwidth}
\centering
\begin{tabular}{|c|c|c|}
\hline
\rowcolor{cyan!30}\multicolumn{3}{|c|}{Studente} \\
\hline
\rowcolor{cyan!30} matricola & nome & cognome \\
\hline
100 & Gianni & Morandi \\
101 & Jerry & Cal\`a \\
102 & Pippo & Franco \\
103 & Mara & Venier \\
\hline
\end{tabular}
\end{minipage}%
\begin{minipage}{.45\textwidth}
\centering
\begin{tabular}{|c|c|c|}
\hline
\rowcolor{cyan!30} \multicolumn{3}{|c|}{Esame} \\
\hline
\rowcolor{cyan!30}matricola & voto & corso \\
\hline
100 & 30 & 27035 \\
100 & 30 & 27038 \\
100 & 30 & 27010 \\
102 & 21 & 27038 \\
90 & 25 & 27045 \\
\hline
\end{tabular}
\end{minipage}
\end{table}
\vspace{2em}
\texttt{SELECT st.matricola, cognome, voto, corso\\FROM Studente AS st LEFT JOIN Esame AS es ON st.matricola=es.matricola;}
\end{frame}
%%%%%%%%%%%%%%%%%
\begin{frame}{LEFT JOIN}
    \begin{table}[h]
\centering
\begin{tabular}{|c|c|c|c|}
\hline
\rowcolor{cyan!30} \multicolumn{4}{|c|}{Result} \\
\hline
\rowcolor{cyan!30} matricola & cognome & voto & corso \\
\hline
100 & Morandi & 30 & 27035 \\
100 & Morandi & 30 & 27038 \\
100 & Morandi & 30 & 27010 \\
101 & Cal\`a & NULL & NULL \\
102 & Franco & 21 & 27038 \\
103 & Venier & NULL & NULL \\
\hline
\end{tabular}
\end{table}
\end{frame}
%%%%%%%%%%%%%%%%%
\begin{frame}{RIGHT JOIN}
\begin{center}
\begin{venndiagram2sets}[shade=red,tikzoptions={opacity=0.5},labelA={},labelB={}, showframe=false, ]
    \fillB
    \setpostvennhook
    {
        \draw (labelA) node[above=2pt] {A};
        \draw (labelB) node[above=2pt] {B};
    }
\end{venndiagram2sets}    
    \end{center}
\end{frame}
%%%%%%%%%%%%%%%%%
\begin{frame}{RIGHT JOIN}
        
\begin{table}[h]
\centering
\begin{minipage}{.45\textwidth}
\centering
\begin{tabular}{|c|c|c|}
\hline
\rowcolor{cyan!30}\multicolumn{3}{|c|}{Studente} \\
\hline
\rowcolor{cyan!30} matricola & nome & cognome \\
\hline
100 & Gianni & Morandi \\
101 & Jerry & Cal\`a \\
102 & Pippo & Franco \\
103 & Mara & Venier \\
\hline
\end{tabular}
\end{minipage}%
\begin{minipage}{.45\textwidth}
\centering
\begin{tabular}{|c|c|c|}
\hline
\rowcolor{cyan!30} \multicolumn{3}{|c|}{Esame} \\
\hline
\rowcolor{cyan!30}matricola & voto & corso \\
\hline
100 & 30 & 27035 \\
100 & 30 & 27038 \\
100 & 30 & 27010 \\
102 & 21 & 27038 \\
90 & 25 & 27045 \\
\hline
\end{tabular}
\end{minipage}
\end{table}
\vspace{2em}
\texttt{SELECT st.matricola, cognome, voto, corso\\FROM Studente AS st RIGHT JOIN Esame AS es ON st.matricola=es.matricola;}
\end{frame}
%%%%%%%%%%%%%%%%%
\begin{frame}{RIGHT JOIN}
    \begin{table}[h]
\centering
\begin{tabular}{|c|c|c|c|}
\hline
\rowcolor{cyan!30} \multicolumn{4}{|c|}{Result} \\
\hline
\rowcolor{cyan!30} matricola & cognome & voto & corso \\
\hline
100 & Morandi & 30 & 27035 \\
100 & Morandi & 30 & 27038 \\
100 & Morandi & 30 & 27010 \\
102 & Franco & 21 & 27038 \\
90 & NULL & 25 & 27045 \\
\hline
\end{tabular}
\end{table}
\end{frame}
%%%%%%%%%%%%%%%%%
\begin{frame}{FULL OUTER JOIN}
\begin{center}
\begin{venndiagram2sets}[shade=red,tikzoptions={opacity=0.5},labelA={},labelB={}, showframe=false, ]
    \fillA
    \fillB
    \setpostvennhook
    {
        \draw (labelA) node[above=2pt] {A};
        \draw (labelB) node[above=2pt] {B};
    }
\end{venndiagram2sets}    
    \end{center}
\end{frame}
%%%%%%%%%%%%%%%%
\begin{frame}{FULL OUTER JOIN}
\begin{table}[h]
\centering
\begin{minipage}{.45\textwidth}
\centering
\begin{tabular}{|c|c|c|}
\hline
\rowcolor{cyan!30}\multicolumn{3}{|c|}{Studente} \\
\hline
\rowcolor{cyan!30} matricola & nome & cognome \\
\hline
100 & Gianni & Morandi \\
101 & Jerry & Cal\`a \\
102 & Pippo & Franco \\
103 & Mara & Venier \\
\hline
\end{tabular}
\end{minipage}%
\begin{minipage}{.45\textwidth}
\centering
\begin{tabular}{|c|c|c|}
\hline
\rowcolor{cyan!30} \multicolumn{3}{|c|}{Esame} \\
\hline
\rowcolor{cyan!30}matricola & voto & corso \\
\hline
100 & 30 & 27035 \\
100 & 30 & 27038 \\
100 & 30 & 27010 \\
102 & 21 & 27038 \\
90 & 25 & 27045 \\
\hline
\end{tabular}
\end{minipage}
\end{table}
\vspace{2em}
\texttt{SELECT st.matricola, cognome, voto, corso\\FROM Studente AS st FULL OUTER JOIN Esame AS es
ON st.matricola=es.matricola;}
\end{frame}
%%%%%%%%%%%%%%%%
\begin{frame}{FULL OUTER JOIN}
\begin{table}[h]
\centering
\begin{tabular}{|c|c|c|c|}
\hline
\rowcolor{cyan!30} \multicolumn{4}{|c|}{Result} \\
\hline
\rowcolor{cyan!30} matricola & cognome & voto & corso \\
\hline
100 & Morandi & 30 & 27035 \\
100 & Morandi & 30 & 27038 \\
100 & Morandi & 30 & 27010 \\
101 & Cal\`a & NULL & NULL \\
102 & Franco & 21 & 27038 \\
103 & Venier & NULL & NULL \\
90 & NULL & 25 & 27045 \\
\hline
\end{tabular}
\end{table}
\end{frame}   
%
\section{Conclusioni}

\begin{frame}{Domande?}
    \begin{figure}
\centering
    \includegraphics[width=0.75\textwidth]{../img/questions.jpg}
\end{figure}
\end{frame}

\begin{frame}{Fine}
    \centering
    \huge Grazie dell'attenzione!
\end{frame}

\end{document}